\setcounter{section}{-1}

\maketitle
\tableofcontents

\section{Introduction}

Welcome to the course notes for Mathematics 1Q1! For a general introduction to the course, including details of assessments, please visit the Course Unit Information module on Canvas. 

\subsection{Help and support}

\textbf{Me:} Carl Whitfield (you can call me Carl).\par
E-mail: \href{mailto:carl.whitfield@manchester.ac.uk} (Start with ``Hi Carl''). \par
Office hour: Mondays 10am-11am, Room 1.214 Alan Turing Building

\bigskip
\textbf{Piazza:} Ask questions (anonymously if you wish) and discuss topics with other students on the course in Piazza, accessed via the Canvas page for this course.

\bigskip
\textbf{Student Support Hub} at the Alan Turing Building:\par
Drop in for Advice, Guidance and Wellbeing Support, Monday to Friday 9:00-17:00

\bigskip
\textbf{School of Natural Sciences Wellbeing Team}: \href{mailto:sons.wellbeing@manchester.ac.uk}{sons.wellbeing@manchester.ac.uk}

\bigskip
\href{https://www.studentsupport.manchester.ac.uk/}{\textbf{University Student Support}}: Central university support and links to the counselling
service.

\bigskip
\href{https://www.dass.manchester.ac.uk/}{\textbf{Disability Advisory and Support Service}}: Register online or book an appointment to
get a University Support Plan.

\subsection{About the course}

The aim of this course provide a comprehensive introduction to concepts in algebra and calculus for year 1 Chemistry students. Below you will find the Intended Learning Outcomes (ILOs) for this course unit.

\begin{ILOs}
    \item \textbf{[ILO1]} Use equations in implicit or parametric form including straight lines or conic sections to extract relevant information.
    \item \textbf{[ILO2]} Identify the appropriate partial fractions form and find the values of any coefficients for a rational function of a single variable.
    \item \textbf{[ILO3]} Differentiate functions and hence locate and classify stationary points for a function of a single variable.
    \item \textbf{[ILO4]} Recall and apply appropriate rules to integrate functions use these to measure areas associated with curves or combinations of curves.
    \item \textbf{[ILO5]} Calculate sums, differences and products involving vectors and use these to solve geometrical problems.
    \item \textbf{[ILO6]} Carry out arithmetic using complex numbers and hence solve polynomial and power equations.
    \item \textbf{[ILO7]} Model chemical processes using differential equation systems and solve these systems using a variety of techniques.
\end{ILOs}

\subsection{How to use these notes}

These notes come in two forms, the ``gappy version''\ifcsname insertBlanks\endcsname (that you are reading now)\fi, which contains gaps in the exercise solutions for you to fill as we work through the material in lectures, and the complete version\ifcsname insertBlanks\endcsname\else(that you are reading now)\fi
, which contains all the full solutions for examples in the lectures and supplementary videos (available on Canvas once the material has been covered in lecture). 

\iflatexml
    You can alter the font size and colour scheme of this html version using the buttons in the top-left and you can download a pdf of these notes by clicking the save button, also in the top-left. 
    \ifcsname insertBlanks\endcsname
        Note that, this html version cannot be annotated directly, if you wish to fill the notes during lectures, you should download the pdf version.
    \fi
\else
    In this pdf version of the notes, some of the html functionality is lost (such as interactive visualisations), to use these, please see the html version in Canvas. 
\fi